% GENERATED FILE. Edit canonical lessons, then run npm run generate:books.
% canonical-source-hash: fnv1a64:bc03dbc251689e0f
% canonical-lessons: MR-C01-dhanyavad, MR-W02-visarga, MR-W02-aa-independent, MR-W02-bha, MR-W02-e-matra, MR-W02-anusvara, MR-W02-ta

\chapter{Hear Thanks, Then Meet Six Signs}
\label{ch:thanks-script-a}
\hlchaptermodality{2}

\begin{chapteropening}
\textbf{By the end of this chapter:} \emph{I can understand and say dhanyavad before decoding it, then write six new Devanagari signs one at a time through visible, guided, and delayed work.}

Write dental ta after a delay while keeping the first five signs separately retrievable.
\end{chapteropening}

\section[dhanyavad]{\mr{धन्यवाद} (dhanyavād) — \textquotedblleft{}thank you\textquotedblright{}}

\label{lesson:MR-C01-dhanyavad}



\begin{quote}
The formal \textquotedblleft{}thanks\textquotedblright{} — Sanskrit, shared across the Indian languages.
\end{quote}

\subsection*{You'll want to know — hear the word before reading it}
Keep the spelling in the heading as a sign you are not yet asked to decode. Listen first: \emph{dhan-ya-vād}. Say it once for \textquotedblleft{}thank you.\textquotedblright{} The next two tiny chapters will give you one missing sign at a time; only then will you read and write the whole word.

\begin{cousinweb}
\emph{Dhanyavād} = Sanskrit \emph{dhanya} (\textquotedblleft{}worthy\textquotedblright{}) + \emph{vāda} (\textquotedblleft{}a saying\textquotedblright{}) — \textquotedblleft{}an utterance of '{[}you are{]} worthy.'\textquotedblright{} The same Sanskrit word behind Hindi \emph{dhanyavād}, Kannada \emph{dhanyavāda}, Telugu \emph{dhanyavādamulu}. Warmer, more spoken Marathi often says \emph{ābhārī āhe} (\textquotedblleft{}{[}I{]} am grateful\textquotedblright{}), on the root \emph{ābhār}, \textquotedblleft{}gratitude.\textquotedblright{}
\end{cousinweb}

\subsection*{Your turn}
Say these aloud:

\begin{itemize}
\raggedright
  \item \textquotedblleft{}dhanyavād\textquotedblright{} — no reading yet
  \item the warmer spoken form — \textquotedblleft{}ābhārī āhe\textquotedblright{}
\end{itemize}

\begin{tcolorbox}[breakable,colback=teal!4,colframe=teal!35!black,title={Before you move on}]
What are the two Sanskrit pieces of \emph{dhanyavād}, and a warmer Marathi alternative? (\emph{dhanya} \textquotedblleft{}worthy\textquotedblright{} + \emph{vāda} \textquotedblleft{}a saying\textquotedblright{}; \emph{ābhārī āhe}.)
\end{tcolorbox}
\section[ḥ]{\mr{ः} — a small breath after a vowel}

\label{lesson:MR-W02-visarga}



\begin{quote}
You already know the greeting \emph{namaskār} and its older first piece, \emph{namaḥ}, \textquotedblleft{}a bow.\textquotedblright{} The final breath in \emph{namaḥ} has its own mark.
\end{quote}

\subsection*{Script}
\begin{quote}
\mr{ः}
\end{quote}

The \textbf{visarga} is two dots written vertically to the right of a syllable. It marks a quiet breath after the vowel. For this first encounter, place the upper dot, lift, then place the lower dot. Stop after this one sign.

\subsection*{Writing — visible model}
Keep \textbf{\mr{ः}} visible. Point to the upper dot, then the lower dot. Finger-trace the two positions once and copy the sign once.

\begin{tcolorbox}[breakable,colback=teal!4,colframe=teal!35!black,title={Before you move on}]
Name it: \textbf{visarga}, a breath after a vowel.

Source: \href{https://www.unicode.org/charts/PDF/U0900.pdf}{Unicode Devanagari chart}, U+0903.
\end{tcolorbox}
\section[ā]{\mr{आ} — long aa standing alone}

\label{lesson:MR-W02-aa-independent}



\begin{quote}
Set one blank line aside for a single new vowel sign.
\end{quote}

\subsection*{Script}
\begin{quote}
\mr{आ}
\end{quote}

This is independent \textbf{ā}, used when the vowel begins a syllable. Curve around the joined left body; lift for the middle shoulder; descend the inner stem; descend the trailing stem; finish with the headline from left to right.

\subsection*{Writing — trace, then guided copy}
Keep \textbf{\mr{आ}} visible. Finger-trace the five runs once. Copy it once, looking back after every run. Say \emph{ā} only when the pen stops.

\begin{tcolorbox}[breakable,colback=teal!4,colframe=teal!35!black,title={Before you move on}]
Is this a vowel sign attached to a consonant? (\textbf{No.} It stands alone.)

Source: \href{https://commons.wikimedia.org/wiki/File:Devanagari\_\%E0\%A4\%86\_stroke\_order.svg}{Saurmandal, Devanagari \mr{आ} stroke order, frames 1–5}.
\end{tcolorbox}
\section[bha]{\mr{भ} — bha, with a puff of breath}

\label{lesson:MR-W02-bha}



\begin{quote}
Make room for one consonant and one careful copy.
\end{quote}

\subsection*{Script}
\begin{quote}
\mr{भ}
\end{quote}

Say \emph{ba}, then \emph{bha} with a small puff against your palm. For the shape, sweep the upper loop into the trunk, lower bowl, and crossbar without lifting; descend the right stem; finish the headline left to right.

\subsection*{Writing — guided copy}
Keep \textbf{\mr{भ}} visible. Copy one run at a time, looking back before each. Say \emph{bha} once after the headline is complete.

\begin{tcolorbox}[breakable,colback=teal!4,colframe=teal!35!black,title={Before you move on}]
Which has the puff: \emph{ba} or \emph{bha}? (\textbf{bha}.)

Source: \href{https://commons.wikimedia.org/wiki/File:Devanagari\_b\%CA\%B0\_\%E0\%A4\%AD.gif}{JackPotte, Devanagari bh stroke animation}.
\end{tcolorbox}
\section[e]{\mr{े} — e above the headline}

\label{lesson:MR-W02-e-matra}



\begin{quote}
Leave a small space above the next practice headline.
\end{quote}

\subsection*{Script}
\begin{quote}
\mr{े}
\end{quote}

This mātrā replaces the inherent vowel with \textbf{e}. Add its short rising stroke above the consonant's headline; it cannot stand alone as a spoken vowel here.

\subsection*{Writing — guided copy and first retrieval}
Copy \textbf{\mr{े}} once beside the visible model. Then cover the previous three models and write, from their names, visarga, independent \emph{ā}, and \emph{bha}. Uncover and repair only one differing part.

\begin{tcolorbox}[breakable,colback=teal!4,colframe=teal!35!black,title={Before you move on}]
Where does \textbf{e} sit? (\textbf{Above the headline.})

Source: \href{https://www.unicode.org/charts/PDF/U0900.pdf}{Unicode Devanagari chart}, U+0947.
\end{tcolorbox}
\section[ṃ]{\mr{ं} — one nasal dot}

\label{lesson:MR-W02-anusvara}



\begin{quote}
Prepare to place one mark above a headline.
\end{quote}

\subsection*{Script}
\begin{quote}
\mr{ं}
\end{quote}

The \textbf{anusvāra} is one dot above the headline. It marks a nasal sound whose exact place follows the sound after it. Place one dot; do not turn it into the two-dot visarga.

\subsection*{Writing — guided copy}
Keep \textbf{\mr{ं}} visible and copy it once. Point above the imaginary headline where the dot belongs.

\begin{tcolorbox}[breakable,colback=teal!4,colframe=teal!35!black,title={Before you move on}]
One dot or two? (\textbf{One.})

Source: \href{https://www.unicode.org/charts/PDF/U0900.pdf}{Unicode Devanagari chart}, U+0902.
\end{tcolorbox}
\section[ta]{\mr{त} — dental ta}

\label{lesson:MR-W02-ta}



\begin{quote}
Touch the tongue lightly behind the upper teeth once.
\end{quote}

\subsection*{Script}
\begin{quote}
\mr{त}
\end{quote}

Touch the tongue lightly behind the upper teeth: \emph{ta}. Sweep left across the shoulder and curve down to the open tip; descend the right stem; finish the headline left to right.

\subsection*{Writing — look, hide, write}
Look at \textbf{\mr{त}} for five seconds. Cover it, wait five seconds, and write it once. Uncover and compare body, stem, then headline.

\begin{tcolorbox}[breakable,colback=teal!4,colframe=teal!35!black,title={Before you move on}]
Say \emph{ta} once with the tongue at the teeth.

Source: \href{https://commons.wikimedia.org/wiki/File:Deva-\%E0\%A4\%A4-order.gif}{Opiaterein, Devanagari ta stroke animation}.
\end{tcolorbox}
